\section{Introduction}\label{sec:intro}
Leakage current is a large part of the power budget of nanometre CMOS circuits and it varies by orders of magnitude with temperature, supply voltage and process parameters such as threshold voltage and oxide thickness~\cite{roy2003}.
Estimating it over many PVT corners with transistor-level SPICE simulation is accurate but slow, because every circuit, input vector and corner needs its own simulation.
A surrogate that predicts the currents of a single transistor from its bias, geometry, temperature and process parameters can be evaluated in microseconds and reused for every circuit built from the same transistors.

This work builds such a hierarchy for the PTM 22\,nm high-performance (HP) model card~\cite{ptm,ptmweb} (BSIM4~\cite{bsim4}) at a nominal supply of 1\,V: \emph{transistor} $\rightarrow$ \emph{gate} $\rightarrow$ \emph{circuit} (Fig.~\ref{fig:flow}).
Four machine-learning models predict the terminal currents $(i_d,i_g,i_s,i_b)$ of an NMOS or PMOS transistor in its ON and OFF state. Only these four models are then used to predict the leakage of an inverter, NAND2, NOR2, AND2 and OR2 for every input combination at 200 random corners, and the leakage of a 3-bit carry-lookahead adder as the sum of its 33 gate leakages; every prediction is compared with a transistor-level SPICE simulation.
The target is a mean accuracy of up to 99\,\%.

\begin{figure}[t]
\centering
\begin{tikzpicture}[font=\footnotesize,
  box/.style={draw,rounded corners=1pt,align=center,minimum width=4.3cm,inner sep=2pt},
  ref/.style={draw,rounded corners=1pt,align=center,fill=gray!12,minimum width=1.9cm,inner sep=2pt},
  >={Stealth[length=1.6mm]}]
\node[box] (A) {Part A: single-transistor SPICE\\ $4\times 10{,}000$ samples};
\node[box,below=4mm of A] (M) {Four models: NMOS/PMOS $\times$ ON/OFF\\ $(V,T,W,P)\rightarrow(i_d,i_g,i_s,i_b)$};
\node[box,below=4mm of M] (G) {Gate leakage, 18 cases $\times$ 200 corners\\ super-node KCL on predicted currents};
\node[box,below=4mm of G] (C) {3-bit CLA, 64 vectors $\times$ 200 corners\\ plain sum of 33 gate leakages};
\node[ref,right=5mm of G] (GS) {Part B\\ SPICE gates};
\node[ref,right=5mm of C] (CS) {Part C\\ SPICE CLA};
\draw[->] (A)--(M); \draw[->] (M)--(G); \draw[->] (G)--(C);
\draw[<->,dashed] (G)--(GS); \draw[<->,dashed] (C)--(CS);
\end{tikzpicture}
\caption{Transistor $\rightarrow$ gate $\rightarrow$ circuit pipeline. The gray boxes are the SPICE references (evaluation only; never used for training or model selection); dashed arrows are the comparisons.}
\label{fig:flow}
\end{figure}

\textbf{Contributions.}
(i)~A reproducible data-generation flow in which all random draws are made inside ngspice, verified against direct model-card edits (Sec.~\ref{sec:data}).
(ii)~Four transistor models chosen on validation data only from four model types, including a tuning step with a pre-declared selection rule (Sec.~\ref{sec:models}).
(iii)~A Kirchhoff (super-node) construction of the gate leakage from predicted terminal currents, which makes the prediction exact for exact currents, whereas the plain current sum fails (Sec.~\ref{sec:gates}).
(iv)~Gate- and circuit-level accuracy together with a ceiling analysis that separates the error of the models from the error of summing isolated gates (Sec.~\ref{sec:circuit}).